\section{Experimental Methodology}
\label{sec:method}

\subsection{Area, Power, Energy, Scalability}
\label{sec:method:ppa}
\subsubsection{RTL Implementation}
Matched \textsc{Tessera} and Planaria~\cite{ghodrati2020planaria} chips were implemented
in RTL at every granularity tier, with the same PE count, granularity
$d$, and on-chip SRAM: four $32{\times}32$ tiers ($1024$ PEs,
$d{=}16$ to $2$). The $16$-sub-array tier reproduces the Planaria
paper's configuration. All designs are
BF16 input with FP32 accumulate and include the strip, sub-array, or
pod memory (one per 4 sub-arrays), the main buffer, and the off-chip
channel models, accounted separately. The classical WS and
unpartitioned TTSA
monoliths calibrated the array level. The Planaria reproduction
implements the full partitioning logic of \S\ref{sec:bg:taxes}
following the released code~\cite{ghodrati2020planaria}.

\subsubsection{Synthesis and Power}
All designs are synthesized with Cadence Genus~\cite{cadence_genus} at
$700$\,MHz in a TSMC 28nm process. We use TSMC 28nm memory compiler to generate SRAM macros.
% The BF16 FMA core was synthesized once and locked as a macro, and every On-chip SRAM is a TSMC N28 memory-compiler macro with its compiler-generated timing and power views.
Each $32$-scale chip ran a
$128{\times}32{\times}32$ GEMM in gate-level simulation, and
Cadence Joules~\cite{cadence_joules} integrated the per-cycle
timeline.



\figreal[t]{eval-ppa.pdf}{eval-area}{1.0}{%
The $32$-scale synthesis, normalized to the
WS-monolith array: (left)~area and (right)~power.}{%
Array-only and full-chip ($+M$ SRAM) totals.}

\figreal[t]{power-timeline.pdf}{power-timeline}{1.0}{%
Array-only dynamic power over one GEMM pass for
WS, Planaria, and \textsc{Tessera}.}{}

\figreal{eval-grainfloor.pdf}{grainfloor}{1.0}{%
Geomean EDP over each layer's own optimum as the minimum partition
granularity coarsens, at HBM4.}{%
Each bar is the geomean over all $259$ layers; the arrow bar spans the
per-benchmark geomeans. Chosen $d{=}8$ in blue.}

\figreal[t]{eval-suite.pdf}{eval-suite}{1.0}{%
Per-network speed-up over the monolithic WS systolic array with HBM4.}{}



\subsection{Cycle-Level Model and Workloads}
\label{sec:method:cycle}
\subsubsection{Cycle-Level Model}
The DNN and multi-tenant experiments use
SCALE-Sim-derived models~\cite{samajdar2020systematic}, extended with the
twist-torus fold timing of \S\ref{sec:ttsa:anchors} and
Algorithm~\ref{alg:tessera-sched}. The twelve-cohort main comparison
uses per-GEMM array timing and a DRAM-cadence model at $1$\,GHz with
HBM4 bandwidth ($2.8$\,TB/s)~\cite{micronhbm}: one logical A/B load and final C store,
without an SRAM capacity limit. Each design uses its own configuration
rule for each stored $(M,N,K)$. With invocation count $r_i$, we report
\begin{equation}
 T_w=\sum_i r_i T_i,\qquad E_w=\sum_i r_i E_i,\qquad
 \mathrm{EDP}_w=E_w T_w.
 \label{eq:workload-service}
\end{equation}
Speed-up is $T_{\rm WS}/T_w$. This measures GEMM service cycles, excluding
arrival-time idle, cross-GEMM packing, and non-GEMM operators. The EDP
ablations below use a separate capacity-aware model.

\subsubsection{Compared Designs}
The designs ran at matched PE count ($128{\times}128$ PEs),
differing in partitioning granularity, dataflow,
and memory organization:
the monolithic classical WS systolic array; Planaria at its published
sixteen $32{\times}32$ sub-arrays;
SOSA~\cite{yuzuguler2022scale}, sixteen partitioned
$32{\times}32$ sub-arrays; FlexSA~\cite{lym2020flexsa}, four $64{\times}64$
quadrants; SISA~\cite{altamura2026sisa}, eight output-stationary
$128{\times}16$ slabs; and \textsc{Tessera} at minimum granularities
$32$ and $8$, denoted \textsc{Tessera}-128-32 and
\textsc{Tessera}-128-8. The latter uses sixteen $128{\times}8$ strips.
For each GEMM, \textsc{Tessera} chooses the EDP minimum among the
$9$ or $25$ logical $(k,n)$ candidates at its respective granularity,
using compute/SRAM/reduction energy.
Planaria-32 uses the composition selected by its native per-GEMM
HBM4 EDP search; the service-cadence Planaria array model then evaluates
that composition at $1$\,GHz.
WS and SOSA have fixed partitions; FlexSA and SISA retain their
phase and slab-fusion rules. Reported cycles correspond to
these choices rather than a separate minimum-cycle search.
The granularity study of
\S\ref{sec:eval:tradeoff} additionally swept \textsc{Tessera} and Planaria
across granularity scales from $d{=}128$ to $2$.

\subsubsection{Workloads}
We evaluate the Planaria network suite and its multi-tenant QoS
scenarios~\cite{ghodrati2020planaria}, and six
Vidur serving traces~\cite{agrawal2024vidur}:
Llama-2-7B~\cite{touvron2023llama2openfoundation},
Llama-3-8B~\cite{grattafiori2024llama3herdmodels},
Llama-2-70B (TP$=$4), and Phi-2~\cite{javaheripi2023phi2} over
the opening $300$ requests of
the Splitwise~\cite{patel2024splitwise} conversation trace plus Llama-2-7B
over the Splitwise code trace and an arXiv
summarization~\cite{cohan2018discourseawareattentionmodelabstractive}
trace at $1$ query per second (QPS). Their saved prefill/decode GEMMs
retain the simulated batch composition and invocation weights.
Phi-2 supplies head dimension $80$, divisible by neither $32$ nor $128$.

Six additional cohorts cover attention and KV-reuse prefill.
Two FlashInfer input collections~\cite{flashinfertrace} contain GQA
decode with $H_q/H_{kv}\in\{4,5\}$ and multi-query attention (MQA)
decode with $H_q=5$ and $H_{kv}=1$, respectively. These are operator-input
snapshots with actual context lengths. MQA names the recorded
single-KV-head operator; its source may be a GQA parallel shard,
so it does not establish a complete MQA model trace. Within each request, the
grouped-head mapping uses QK $(M,N,K)=(qH_q/H_{kv},L,D_h)$ and PV
$(qH_q/H_{kv},D_h,L)$, where $q$ is the query-token count, $L$ the
context length, and $D_h$ the head dimension. Batch and KV-head counts
enter $r_i$ once; requests with private KV caches remain separate. This
distinguishes grouped attention GEMMs from single-token projection/FFN
GEMVs with $M=1$.

The other four cohorts use CacheBlend SAMSum and WikiMQA~\cite{cacheblendcode},
and EPIC HotpotQA and Multi-News~\cite{epiccode,longbenchdata}.
We tokenize the upstream requests using the pinned
Mistral-7B-Instruct-v0.2 configuration~\cite{mistralv02config} and derive
all decoder GEMMs from each framework's recomputation policy, including
QKV, output projection, fused gate/up, down projection, attention, and
the final one-row vocabulary projection. Each source contributes its
first $200$ requests: $799$ of $800$ are valid; one Multi-News request
has out-of-range upstream gather indices. Its input and failure reason
are retained. These are source-derived prefill plans at batch~1 and TP~1.
The artifact records versions, tokens, exclusions, and repeat counts.
Each cohort is reported separately: snapshots establish operator
shapes, simulated traces supply serving GEMMs, and source-derived plans
supply reuse GEMMs; none supplies measured end-to-end GPU latency.

\subsubsection{EDP Ablation Controls}
The capacity-aware ablation model uses $16{,}384$ PEs at each minimum granularity
$d\in\{32,8\}$. Independent WS arrays number $16$ at $d=32$ and $256$
at $d=8$. The physical-square variant permits unequal logical
$(k,n)$ tiles but pads each physical region to a square; the free
variant includes those candidates and also permits rectangles.
Padding consumes PE area and cycles, not additional useful MACs or
operand accesses. Each variant independently minimizes modeled EDP.

Each GEMM receives the HBM traffic of its EDP-selected native
Planaria-32 configuration: unrestricted, square-composed, and independent
cores for the free, physical-square, and independent variants,
respectively. Each record is fixed across that variant's candidate
search at both granularities. It retains native private-buffer reuse,
padded transfers, initial output reads, and 64-bit off-chip partials.
This external traffic sensitivity does not transplant native thread
scheduling or run Planaria-8. \textsc{Tessera}'s compute/SRAM candidates use
$12{,}000{,}000$~B, whereas the saved native traffic comes from
$12$~MiB. Its on-chip partial sums remain FP32.

For initial, middle, and final transfer bytes, $I,X,F$ are the
separately rounded transfer times at $2800$~B/cycle. With compute time
$C$, the aggregate overlap model is $T=I+\max(C,X)+F$; it does not
validate event-level operand arrival. Energy retains $6$~pJ per operation,
$0.95/2.2$~pJ per input/output SRAM word, and the existing analytic
reduction cost. HBM reads and writes each cost $32$~pJ/B, an inherited
modeling assumption rather than a measured HBM4 value.
SRAM-port bandwidth is unconstrained. These coefficients are separate
from the gate-level measurements in
\S\ref{sec:method:ppa}.

The skew ablation fixes the free variant's selected mapping, packing,
issue interval, bank release, traffic, and energy. It adds only the
exposed skew tail to $C$ before recomputing overlap: $H-1$ for a
uniform $H{\times}W$ compute chunk, or
$\max(H_i+W_i-2)-\max(W_i-1)$ over its last round for mixed shapes.
Intermediate-round tails remain hidden under the model;
this changes the tail term alone, not the complete WS schedule.

% eval figure declarations moved to the \S V end (2026-08-01) so
% Fig 8--11 float onto the eval pages from an early declaration point.
